% !TeX program = tectonic
\documentclass[11pt,a4paper]{article}
\usepackage{fontspec}
\usepackage[russian]{babel}
\babelfont{rm}[Language=Default]{Liberation Serif}
\babelfont[Language=Default]{sf}{Carlito}
\babelfont[Language=Default]{tt}{DejaVu Sans Mono}
\usepackage{amsmath,amssymb,amsthm}
\usepackage{geometry}
\geometry{a4paper, top=2.4cm, bottom=2.4cm, left=2.4cm, right=2.4cm}
\usepackage{booktabs,tabularx,enumitem,xcolor,tcolorbox}
\tcbuselibrary{breakable,skins}
\definecolor{linkblue}{HTML}{1F4E79}
\definecolor{statgreen}{HTML}{2F6B3A}
\definecolor{goldacc}{HTML}{8B7E5A}
\usepackage[colorlinks=true,linkcolor=linkblue,urlcolor=linkblue,unicode=true]{hyperref}
\theoremstyle{plain}
\newtheorem{theorem}{Теорема}
\newtheorem{lemma}[theorem]{Лемма}
\newtheorem{proposition}[theorem]{Предложение}
\newtheorem{corollary}[theorem]{Следствие}
\theoremstyle{definition}
\newtheorem{definition}[theorem]{Определение}
\theoremstyle{remark}
\newtheorem{remark}[theorem]{Замечание}
\newtcolorbox{statusbox}[1][]{colback=statgreen!5,colframe=statgreen!75,
  title={\textbf{Сводка доказанного:} #1},fonttitle=\small,boxrule=0.7pt,arc=2pt,breakable}
\newtcolorbox{databox}[1][]{colback=goldacc!6,colframe=goldacc!80,
  title={\textbf{Данные протокола:} #1},fonttitle=\small,boxrule=0.7pt,arc=2pt,breakable}
\newtcolorbox{verifybox}[1][]{colback=linkblue!5,colframe=linkblue!80,
  title={\textbf{Верификация:} #1},fonttitle=\small,boxrule=0.7pt,arc=2pt,breakable}
\widowpenalty=10000
\clubpenalty=10000
\setlength{\parskip}{0.35em}
\newcommand{\CC}{\mathbb{C}}
\newcommand{\RR}{\mathbb{R}}
\newcommand{\ZZ}{\mathbb{Z}}
\newcommand{\QQ}{\mathbb{Q}}
\newcommand{\Ga}{\Gamma}
\newcommand{\Bfun}{\mathrm{B}}
\newcommand{\im}{\operatorname{Im}}
\newcommand{\re}{\operatorname{Re}}
\newcommand{\lcm}{\operatorname{lcm}}
\newcommand{\SNF}{\operatorname{SNF}}
\newcommand{\Jac}{\operatorname{Jac}}
\newcommand{\rank}{\operatorname{rank}}
\newcommand{\Ind}{\operatorname{Index}}
\newcommand{\disc}{\operatorname{disc}}
\begin{document}
\begin{center}
{\LARGE\bfseries Стенд K3: ранг Нерона–Севери 20}\\[0.4em]
{\large сорок восемь прямых и точная матрица пересечений}\\[0.6em]
{\normalsize Исаев Исхак Хамзатович}\\[0.2em]
{\small Программа <<Динамический принцип>> --- лаборатория \texttt{hodge-laboratory}}\\[0.15em]
{\small Отдельная монография теоремы · Монография 8/16}
\end{center}
\vspace{0.4em}
{\color{linkblue}\hrule height 0.8pt}
\vspace{0.8em}

\section{Постановка}

Фермат-квартика $x^4+y^4+z^4+w^4=0$ в $\mathbb{P}^3$ --- K3-поверхность
максимального ранга Нерона--Севери. Гипотеза Ходжа для неё доказана
(Шиода, 1979); программа воспроизводит все инварианты точной
целочисленной арифметикой над $\QQ$.

\begin{theorem}[инварианты K3-стенда]\label{th:k3stand}
Для точной матрицы пересечений $49\times49$ (48 прямых + класс
гиперплоскости $h$) Фермат-квартики выполнено:
(1) $\rank_\QQ G=20$ --- совпадает с теоремой Шиода $\rho=20$;
(2) сигнатура численной части равна $(1,19)$;
(3) определитель подрешётки равен $64=8^2$;
(4) сумма четырёх прямых одного семейства эквивалентна $h$:
$[L_1(a,0)]+[L_1(a,1)]+[L_1(a,2)]+[L_1(a,3)]\sim h$.
\end{theorem}

\begin{proof}
Прямые: три семейства по $16$:
$L_1(a,b)\colon x+i^ay=0$, $z+i^bw=0$;
$L_2(a,b)\colon x+i^az=0$, $y+i^bw=0$;
$L_3(a,b)\colon x+i^aw=0$, $y+i^bz=0$, $a,b\in\ZZ/4$.
Правила пересечений комбинаторны: $L\cdot L=-2$; внутри семейства
пересечение $1$ iff ровно один индекс совпадает; между семействами
циклические условия ($L_1\cdot L_2$: $a_1+b_2\equiv a_2+b_1$;
$L_2\cdot L_3$: $a_1+b_1\equiv a_2+b_2$; $L_1\cdot L_3$:
$c\equiv a+b+d+1$ --- лишняя единица из суммы трёх нечётных
показателей); $h\cdot L=1$, $h^2=4$.
Ранг: точное целочисленное приведение (фракционная элиминация без
округлений) --- $20$ ненулевых строк; двойная регистрация через
миноры Смит-нормы.
Сигнатура: собственные значения --- одно положительное, $19$
отрицательных, нулей нет (порог отделения $0.1$).
Определитель: полная решётка K3 унимодулярна сигнатуры $(3,19)$;
индекс подрешётки $8$, $\det=64=8^2$.
Эквивалентность: вектор пересечений суммы прямых покомпонентно равен
вектору $h$ (все $49$ компонент); разность ортогональна
алгебраической подрешётке и лежит в $H^{2,0}\oplus H^{0,2}$ по
индексу Ходжа; но разность алгебраична, а алгебраические классы
ортогональны $H^{2,0}$; значит, разность нулевая.
\end{proof}

\begin{proposition}[циклы кодимерсии 2]
Класс точки: $[\mathrm{pt}]=h^2/4$ --- точно, так как
$(h^2)^2=\deg X=4$ и $H^4(X)=\QQ h^2$. На $K3\times K3$: пары
$(c_2,h_2)$ дают $24$ независимых класса кодимерсии $(2,2)$,
диагональные $h^2$-пары --- ещё $4$; все алгебраичны с явными
носителями. Слепой зоны нет: каждый $(2,2)$-класс, ортогональный
построенным циклам, нулевой по индексу Ходжа.
\end{proposition}

\begin{databox}[числа]
$n_{\mathrm{lines}}=48$; $\rank=20=\rho$ (Шиода); инерция $(1,19)$;
$\det=64$; гиперплоскостное соотношение точно; орбит $12$; кратности
характеров $(1,7,7,7)$ --- суммарно $22$; кодимерсия 2: $4+24$;
время счёта $1.04$~с.
\end{databox}

\begin{statusbox}[итог]
Все инварианты K3-стенда вычислены точно над $\QQ$; совпадение с
теоремой Шиода подтверждено; эквивалентность семейства и $h$
доказана точно.
\end{statusbox}

\begin{verifybox}[как воспроизвести]
\texttt{python3 laboratory.py} $\to$ <<Стенды $\to$ K3>>: матрица
$49\times49$, ранг, сигнатура, определитель --- за секунду.
\end{verifybox}

\vfill
{\color{linkblue}\hrule height 0.6pt}
\vspace{0.5em}
{\small\color{gray}
Лицензия: индивидуальная исключительная лицензия Исаева Исхака Хамзатовича.
Все права на вычисления, формулы и тексты принадлежат автору.
Верификация: \texttt{python3 laboratory.py} (пункт <<Стенды>>/<<Сертификаты>>),
Lean: \texttt{verification/lean/HodgeLaboratory.lean}.
}
\end{document}
